% Chapter 5 - the standalone blob.

\guidechapter{OTHER TOOLCHAINS}{%
  \item Building and identifying the blob
  \item The stable jump table
  \item The parameter block
  \item Ultimate Audio and software-vsprite extension entries
  \item Calling it from BASIC
  \item Choosing a different base address
}

\section{Build And Identify The Blob}

The blob is the SDK linked as one binary with a table of jump instructions at
a fixed address. Any assembler or compiler can call it, because calling it
needs no linker and no object format: you load the binary and \cmd{jsr} into
the table.

\cmd{make blob} produces \cmd{bindings/blob/build/ultimate-7000.bin}. Load it
at \cmd{\$7000}. The first three bytes are the PETSCII signature \cmd{UCI};
the fourth is the version of the ABI, the calling interface.

\begin{note}
The blob has to sit below \cmd{\$A000}, because \cmd{\$A000--\$BFFF} is BASIC
ROM. A program can reach every byte of a blob at \cmd{\$7000} with the machine
exactly as it found it. One that ran past \cmd{\$A000} would have to switch
BASIC ROM out and back around every call.
\end{note}

\begin{warning}
\cmd{\$7000} is inside BASIC's memory, not above it. BASIC owns
\cmd{\$0801--\$9FFF}: the program area grows up from \cmd{\$0801} and strings
and arrays grow down from \cmd{\$9FFF}. A program that has replaced BASIC owns
that memory instead and needs to do nothing. A program that leaves BASIC running
must lower the top of BASIC memory to the base address before it loads the blob,
or BASIC will write over it. \emph{Call The Blob From BASIC}, later in this
chapter, shows the two \bas{poke}s and the \bas{clr} that do it.
\end{warning}

This example checks the default build before calling it:

\begin{listingbox}
\hspace*{2em}lda \$7000\\
\hspace*{2em}cmp \#\$D5\hspace{2em}; PETSCII 'U'\\
\hspace*{2em}bne no\_sdk\\
\hspace*{2em}lda \$7001\\
\hspace*{2em}cmp \#\$C3\hspace{2em}; PETSCII 'C'\\
\hspace*{2em}bne no\_sdk\\
\hspace*{2em}lda \$7002\\
\hspace*{2em}cmp \#\$C9\hspace{2em}; PETSCII 'I'\\
\hspace*{2em}bne no\_sdk\\
\hspace*{2em}lda \$7003\\
\hspace*{2em}cmp \#1\hspace{3em}; ABI version\\
\hspace*{2em}bne wrong\_version
\end{listingbox}

\begin{resultbox}
All four comparisons match; execution continues with ABI version 1.
\end{resultbox}

The version byte changes only when an existing entry changes meaning. New
entries are appended, so a program written against version~1 keeps working
when the table grows.

\evidence{\cmd{bindings/blob/Makefile}; \cmd{bindings/blob/blob.s};
\cmd{tests/emulator/blob.suite}}

\section{Call The Stable Jump Table}

Entries begin at base plus \cmd{\$04}. Each entry is a three-byte \cmd{jmp};
new entries are appended so published offsets do not move.

\begin{listingbox}
\hspace*{2em}jsr \$7004\hspace{2em}; +\$04  uci\_init\\
\hspace*{2em}jsr \$701C\hspace{2em}; +\$1C  ultimate\_init
\end{listingbox}

\begin{resultbox}
A = ULTIMATE_OK after each initialization call.
\end{resultbox}

These are the two initialization entries. Use \cmd{ultimate\_init} for
application code; the lower-level \cmd{uci\_init} entry is exposed for
transport users. Both return a result in \cmd{A}.

\subsection*{The two kinds of entry}

The table has two halves, and the difference decides how you read a result.

Entries from \cmd{+\$04} to \cmd{+\$4C} pass their arguments in registers, the
way the assembly chapter's entry points do, and return the result in \cmd{A}
and nowhere else.

Entries from \cmd{+\$4F} upwards take their arguments from the parameter block
described in the next section, because a caller cannot pass a filename in
\cmd{A} and \cmd{X}. Each of these also stores its result in \cmd{bp\_result}
before returning it in \cmd{A}.

That second group is what a BASIC program can drive, because \bas{sys} can
neither pass a register nor read one back. The extension tables at
\cmd{+\$2E8} and \cmd{+\$300} contain both kinds; their individual contracts
say whether they use registers or the parameter block.

The complete offset table is in \cmd{bindings/blob/README.md}. Read the
offsets there rather than counting entries: the table gained the file
information, disk image, clock, machine and HTTP body services after the
offsets printed in this chapter were published.

\evidence{\cmd{bindings/blob/blob.s}; \cmd{bindings/blob/README.md};
\cmd{tests/emulator/blob.suite}}

\Needspace{20\baselineskip}
\section{Use The Parameter Block}

Service entries that need several arguments use a 512-byte block at
base plus \cmd{\$100}. These are the first four fields:

\begin{reftable}{llp{0.48\linewidth}}
Offset & Field & Purpose \\
\midrule
\cmd{+\$00} & \cmd{bp\_result} & last service result \\
\cmd{+\$05} & \cmd{bp\_len} & length in or out \\
\cmd{+\$29} & \cmd{bp\_name} & 40-byte area for a terminated name \\
\cmd{+\$51} & \cmd{bp\_reply} & 256-byte reply area \\
\end{reftable}

Later fields extend beyond the first page while retaining their published
offsets. Consult the blob README for the complete layout rather than deriving
an address from this subset.

\evidence{\cmd{bindings/blob/blob.s}; \cmd{bindings/blob/README.md};
\cmd{tests/emulator/blob.suite}}

\Needspace{34\baselineskip}
\section{Call Audio And Vsprite Extensions}

Ultimate Audio occupies the extension entries from \cmd{+\$2E8} through
\cmd{+\$300}. The normal sequence is
\cmd{ultimate\_audio\_init}, optionally
\cmd{audio\_load\_wav}, then \cmd{audio\_configure} and
\cmd{audio\_start}. Configure, start, stop and clear use the
\cmd{ultimate\_audio\_voice} structure at
\cmd{BLOB + BLOB\_BP\_AUDIO}; the WAV entry also uses
\cmd{BLOB\_BP\_NAME} and \cmd{BLOB\_BP\_REU}. Chapter~3 defines every voice
field and the safe configure/start/stop sequence. The complete stable offset
table is in \cmd{bindings/blob/README.md}.

The software-vsprite entry is \cmd{BLOB + BLOB\_VSPRITE\_DRAW}, currently
\cmd{+\$303}. Put the descriptor address in \cmd{BLOB\_BP\_ADDR} and read the
result from \cmd{BLOB\_BP\_RESULT}:

\begin{listingbox}
BLOB = \$7000\
\
\hspace*{2em}lda \#<sprite\
\hspace*{2em}sta BLOB + BLOB\_BP\_ADDR\
\hspace*{2em}lda \#>sprite\
\hspace*{2em}sta BLOB + BLOB\_BP\_ADDR + 1\
\hspace*{2em}jsr BLOB + BLOB\_VSPRITE\_DRAW\
\hspace*{2em}lda BLOB + BLOB\_BP\_RESULT\
\hspace*{2em}bne draw\_failed\
\
sprite:\
\hspace*{2em}.word \$6000, image, mask, 0\
\hspace*{2em}.byte 12, 72, 2, 8, 0, VSPRITE\_F\_MASKED
\end{listingbox}

\begin{resultbox}
The masked image is composited into the bitmap; the parameter result is zero.
\end{resultbox}

The descriptor and flag contract is the one in \emph{Composite A Software
Vsprite} in Chapter~3. This is a parameter-block shim around the same assembly
routine, not another renderer.

\evidence{\cmd{bindings/blob/blob.s}; \cmd{bindings/blob/README.md};
\cmd{bindings/asm/uci\_protocol.inc}; \cmd{src/uci/vsprite.s}}

\section{Call The Blob From BASIC}

The page-aligned block gives BASIC a calling convention it can express with
\bas{poke}, \bas{peek} and \bas{sys}. Before any of that, BASIC has to be told
to stay out of the memory the blob occupies. Type this first, then load the
blob:

\begin{screen}
POKE 56,112 : POKE 55,0 : CLR
\end{screen}

112 is \cmd{\$70}, the high byte of the base address, and 55 and 56 are the low
and high bytes of BASIC's top-of-memory pointer at \cmd{\$37--\$38}. \bas{clr}
is what makes it take effect, because it resets the string pointer from the new
top. BASIC then keeps \cmd{\$0801--\$6FFF}, and the blob's code at \cmd{\$7000}
and its variables at \cmd{\$9F00} are both out of its reach. Loading the blob
without this leaves it where BASIC's strings and arrays will eventually be
written.

This program then checks the signature, brings the SDK up, confirms that an
Ultimate answered, and prints the size of the RAM expansion:

\begin{screen}
10 B = 28672 : P = B + 256\\
20 IF PEEK(B)<>213 OR PEEK(B+1)<>195 THEN PRINT "NO SDK":END\\
25 IF PEEK(B+2)<>201 THEN PRINT "NO SDK":END\\
30 SYS B + 28\\
40 SYS B + 82 : IF PEEK(P) THEN PRINT "NO ULTIMATE":END\\
50 SYS B + 175\\
60 PRINT "REU BANKS:"; PEEK(P+5) + PEEK(P+6)*256
\end{screen}

Line~10 sets \bas{b} to the blob's base, decimal 28672 for \cmd{\$7000}, and
\bas{p} to the parameter block a page above it. Lines~20 and~25 compare the
three signature bytes; 213, 195 and 201 are \cmd{\$D5}, \cmd{\$C3} and
\cmd{\$C9}.

Line~30 calls \cmd{ultimate\_init} at offset \cmd{\$1C}. That is a register
entry, so it returns its result in \cmd{A} and writes nothing to the block,
and \bas{sys} cannot read \cmd{A}. Line~40 therefore calls \cmd{getpath} at
offset \cmd{\$52}, which is a parameter-block entry: it asks the Ultimate for
the current directory and leaves the result in \cmd{bp\_result}, so
\bas{peek(p)} is the first result this program can test. A non-zero value
there means no Ultimate answered.

Line~50 calls \cmd{reu\_size} at offset \cmd{\$AF}, which puts the 16-bit bank
count in \cmd{bp\_len}. Line~60 reads it back as two bytes, low byte first, and
prints it. A machine with a 2MB expansion prints 32; one with no expansion
prints 0.

\begin{warning}
Do not check \bas{peek(p)} after a register entry such as \cmd{\$1C}. Those
entries never write \cmd{bp\_result}, so the byte still holds whatever the
previous call left there, and in a freshly loaded blob that is zero. The test
would pass whether the call worked or not.
\end{warning}

\evidence{\cmd{bindings/blob/blob.s}; \cmd{bindings/blob/README.md};
\cmd{tests/emulator/blob.suite}}

\section{Choose A Different Base Address}

Prefer building at the final base: \cmd{make -C bindings/blob BASE=6000} links
a \cmd{\$6000} image directly, and \cmd{VARS} and \cmd{ZP} move the SDK's RAM
the same way.

When the address is known only at run time, the build also emits a relocation
table beside the binary for \cmd{reloc.s} to apply. The relocated emulator
suite copies the default blob to \cmd{\$6000}, relocates the copy, calls it,
and then compares it byte for byte against an image the linker built for
\cmd{\$6000} directly.

\evidence{\cmd{bindings/blob/Makefile}; \cmd{bindings/blob/reloc.s};
\cmd{tests/emulator/blob-relocated.suite.in};
\cmd{tests/emulator/relocharness.s}}
